\section{Literature Review}

This section reviews existing methods for automatic chart understanding, organized by
approach paradigm. We analyze the strengths and limitations of each category to motivate
the hybrid design of Geo-SLM.

\subsection{End-to-End Multimodal Approaches}

Recent advances in vision-language models have enabled end-to-end chart understanding
systems that process chart images directly without explicit geometric decomposition.

\textbf{DePlot} \cite{liu2023deplot} employs a two-stage approach: a Pix2Struct model
translates the chart image into a linearized data table, which is then processed by a
large language model (LLM) for question answering. DePlot achieves competitive results
on the ChartQA benchmark \cite{masry2022chartqa} with one-shot prompting. However,
because the Pix2Struct model estimates tabular values from visual patterns, it is prone
to hallucination -- producing plausible but numerically incorrect entries, especially
for charts with dense data points or overlapping series.

\textbf{MatCha} \cite{lee2023matcha} extends the Pix2Struct architecture with two
pretraining objectives: math reasoning and chart derendering. The derendering task trains
the model to generate the underlying data table from a chart image, improving numerical
grounding. MatCha achieves strong results on both ChartQA and PlotQA \cite{methani2020plotqa}
benchmarks. Despite these improvements, the model remains a large (282M+ parameters)
encoder-decoder that requires significant compute resources and does not provide
interpretable intermediate outputs.

\textbf{ChartReader} \cite{chen2022chartreader} proposes a unified framework that
detects chart components (bars, lines, legend entries) using Faster R-CNN, then applies
heuristic rules to reconstruct the data table. This approach provides more interpretable
outputs than DePlot or MatCha, but its component detection relies heavily on annotated
training data and struggles with unconventional chart layouts.

\subsection{Classical Geometric Approaches}

Before deep learning, chart understanding relied on classical image processing and
geometric analysis.

\textbf{ReVision} \cite{savva2011revision} classifies charts using visual features and
extracts data from bar charts using Hough line detection. While fully deterministic and
interpretable, ReVision supports only a single chart type (bar) and requires strict
assumptions about chart layout (aligned bars, visible grid lines). The method fails
on scatter plots, line charts, and charts with non-standard aesthetics.

Other traditional approaches use contour detection, template matching, and projection
profiles to identify chart elements. These methods offer high precision for well-formatted
charts but lack robustness against the visual diversity found in real-world academic
publications.

\subsection{Object Detection for Document Analysis}

The YOLO family \cite{redmon2016yolo, jocher2023yolov8} has been widely adopted for
object detection tasks in document analysis. YOLOv8 provides real-time detection with
competitive accuracy, making it suitable for chart region localization within full-page
document images. For chart detection specifically, the single-class detection strategy
(detecting ``chart'' as one class) with a separate classification stage has been shown
to outperform multi-class detection, since YOLO excels at localization while dedicated
classifiers (e.g., ResNet-18 \cite{he2016resnet}) provide higher type-classification
accuracy.

\subsection{OCR and Text Extraction}

Optical character recognition is fundamental to extracting textual information from charts,
including axis labels, titles, legend entries, and tick mark values. PaddleOCR
\cite{du2020paddleocr} and EasyOCR \cite{jaided2020easyocr} are two widely-used
open-source frameworks supporting multi-language text detection and recognition.
In the context of chart analysis, OCR faces unique challenges: text may be rotated
(Y-axis labels), overlapping (dense tick marks), or embedded as vector graphics
(invisible to pixel-based OCR). Post-processing heuristics -- such as spatial role
classification (assigning text to title, axis, legend, or value roles based on position)
and common-error correction (e.g., ``O'' to ``0'', ``l'' to ``1'') -- are essential
for reliable chart text extraction.

\subsection{Small Language Models and Parameter-Efficient Fine-Tuning}

Large language models have demonstrated strong reasoning capabilities for chart
understanding, but their deployment on consumer hardware is impractical. Small language
models (SLMs) with 1--3 billion parameters offer a practical alternative. Qwen-2.5
\cite{yang2024qwen2} provides instruction-tuned models starting at 0.5B parameters,
with the 1.5B variant balancing capability and efficiency.

Parameter-efficient fine-tuning methods allow adapting these models to domain-specific
tasks without full retraining. LoRA \cite{hu2022lora} introduces low-rank decomposition
of weight update matrices, training only 0.1--1\% of model parameters. QLoRA
\cite{dettmers2024qlora} extends LoRA with 4-bit quantization (NF4 format), reducing
memory requirements to as low as 4~GB VRAM for a 1.5B-parameter model. This enables
fine-tuning and inference on consumer GPUs such as the NVIDIA RTX 3060 (6~GB VRAM).

\subsection{Summary and Research Gap}

\Cref{tab:related_work} summarizes the comparison of existing approaches.

\begin{table}[htbp]
\centering
\caption{Comparison of existing chart understanding methods.}
\label{tab:related_work}
\begin{tabular}{@{}lp{2.5cm}rrlc@{}}
\toprule
\textbf{Method} & \textbf{Approach} & \textbf{Value Acc.} & \textbf{Types} & \textbf{Model Size} & \textbf{Explainable} \\
\midrule
DePlot \cite{liu2023deplot}     & Pix2Struct + LLM       & $\sim$85\% & 5+ & 282M+ & No \\
MatCha \cite{lee2023matcha}     & Chart pretraining      & $\sim$80\% & 5+ & 282M+ & No \\
ChartReader \cite{chen2022chartreader} & Detection + rules & $\sim$75\% & 3  & N/A   & Partial \\
ReVision \cite{savva2011revision}      & Hough transform   & $\sim$70\% & 1  & N/A   & Yes \\
\midrule
\textbf{Geo-SLM (Ours)} & \textbf{Hybrid neuro-symbolic} & \textbf{Target $>$95\%} & \textbf{8} & \textbf{1.5B} & \textbf{Yes} \\
\bottomrule
\end{tabular}
\end{table}

The key research gap is the absence of a system that combines:
\begin{enumerate}
    \item \textbf{Neural perception} for robust detection and classification across chart types,
    \item \textbf{Geometric measurement} for deterministic, traceable value extraction, and
    \item \textbf{Local SLM reasoning} for semantic understanding without cloud dependency.
\end{enumerate}

Geo-SLM addresses this gap by integrating all three paradigms in a modular, 5-stage
pipeline where each component operates on well-defined schema contracts, enabling
independent development, testing, and replacement of any component.

